\documentclass[../main.tex]{subfiles}
\begin{document}
\section{連続読み出し式データ収集システム}
\label{sec:background}
なぜ従来のトリガー型DAQシステムではなく連続読み出し式DAQを用いるか述べる

連続読み出し式DAQシステムの違いに触れつつ、連続読み出し式DAQシステムとはなにか、どんな原理か、どんなFEEを用いるかについてまとめる
\subsection{トリガー型DAQシステム}
\subsection{連続読み出し式DAQシステム}
フロントエンド回路が外部のトリガー信号を受けてヒット信号を転送するのではなく、
\subsection{MARQに用いる連続読み出し式DAQシステム}
連続読み出し式DAQ+オンラインフィルタ on it
\subsubsection{ハードウェア}
どんなFEEを使うか
\begin{tcolorbox}[
  title=執筆メモ,
  colback=yellow!8,
  colframe=black!30,
  breakable
]
\begin{itemize}
  \item AMANEQ
  \item CIRASAME
  \item ASAGI
\end{itemize}
\end{tcolorbox}

\subsubsection{ソフトウェア}
どんなソフトウェア(フレームワーク)を使うか
\begin{tcolorbox}[
  title=執筆メモ,
  colback=yellow!8,
  colframe=black!30,
  breakable
]
\begin{itemize}
  \item NestDAQ
  \begin{itemize}
    \item 連続読み出し式DAQを実現するためのソフトウェアフレームワーク
    \item 分散型DAQシステムを管理する
    \begin{itemize}
      \item 複数プロセス
      \item 複数ノード
    \end{itemize}
    \item いくつかのライブラリを使い、以下の機能を提供する
    \begin{itemize}
      \item プロセスの状態管理(プロセスが状態機械を持つ)
      \item 通信チャンネル(プロセス間のメッセージ通信機能)
      \item トポロジーの生成, サービス発見
      \item RUN制御
      \item 設定管理(パラメータの管理)
      \item モニタリング(各デバイスの動作状態, データ転送量)
    \end{itemize}
  \end{itemize}
  \item nestdaq-user-impl
  \begin{itemize}
    \item NestDAQ上で動作するユーザータスク実装群
    \item FEEからデータを受け取ってからデータを外部に出力するまでの処理を定義する
    \begin{itemize}
      \item ディスクに保存
      \item 間引いてディスクに保存
      \item イベントディスプレイに表示
    \end{itemize}
    \item つまりオンライン処理
    \item どんな実験にも共通して使われる処理を行うためのNestDAQデバイスはすでに(\url{https://github.com/spadi-alliance/nestdaq-user-impl.git})で提供されている
    \item NestDAQデバイスをどのように組み合わせるかをトポロジーが定義する
    \item 最小構成のDAQならFEEからデータを受け取ってディスクに書き込むような、Sampler→FileSink
    \item オンラインでデータ削減を行うための仕組みが提供されている
    \begin{itemize}
      \item Sampler→STFBuilder→TFBuilder→LogicFilter→TimeFrameSlicer→HighLevelFilter→FileSink
    \end{itemize}
  \end{itemize}
\end{itemize}
\end{tcolorbox}
\subsection{NestDAQ上に実装できるオンラインフィルタ}
\end{document}
